%! Piecewise-Defined Functions — Unit 2, Lesson 2.6 — Slides (Beamer)
\documentclass[aspectratio=169, 11pt]{beamer}
\usepackage{atda-beamer}   % provides \royalheader{} and \sectionlabel[color]{}
\usepackage{pgfplots}      % atda-beamer does NOT load pgfplots
\pgfplotsset{compat=1.18}

\begin{document}

% TITLE SLIDE (CourseName is not defined in beamer, so the name is literal)
{
\setbeamercolor{background canvas}{bg=royal}
\begin{frame}[plain]
  \vspace{2.8cm}\hspace{0.55cm}%
  \begin{minipage}{0.9\textwidth}
    {\color{goldacc}\bfseries\small LESSON 2.6}\\[0.5cm]
    {\color{white}\bfseries\LARGE Piecewise-Defined Functions}\\[1.2cm]
    {\color{mistmid}\small Algebra, Trigonometry, and Data Analysis~$\cdot$~Unit 2}
  \end{minipage}
\end{frame}
}

% ── HOOK ──────────────────────────────────────────────────────────────────────
\begin{frame}[t, plain]
  \royalheader{Two Students, One Paycheck, and \$60 Nobody Can Find}
  \sectionlabel[goldacc]{hook --- \$12 an hour for the first $40$ hours, \$18 an hour after that}
  A student worked \textbf{$45$ hours} last week. What is the pay?
  \begin{block}{Both used both rates, and neither made an arithmetic mistake}
    \small
    \textbf{A:} ``$45$ hours at \$12 is \$540, plus $5$ overtime hours at \$18 is \$90. \textbf{So
    \$630.}''\\[0.15cm]
    \textbf{B:} ``$40$ hours at \$12 is \$480, plus $5$ hours at \$18 is \$90. \textbf{So \$570.}''
  \end{block}
  \vspace{0.1cm}
  \textbf{So the only question left is: which hours does each rate apply to?}
\end{frame}

% ── THE RESOLUTION ────────────────────────────────────────────────────────────
\begin{frame}[t, plain]
  \royalheader{A Rule Owns Only Its Own Stretch of Inputs}
  \sectionlabel{that sentence is the entire content of the word ``piecewise''}
  \begin{block}{Student B answered it}
    \centering\large
    \textbf{Hours $41$--$45$ are paid at \$18 \emph{instead of} \$12 --- not on top of it.}
  \end{block}
  \begin{itemize}
    \setlength{\itemsep}{0.3cm}
    \item Student A charged the \$12 rate to \emph{all} $45$ hours, so those $5$ hours got paid
          twice: $5 \times \$12 = \$60$ too much.
    \item Before you use a rule, ask: \textbf{which piece owns this input?}
    \item That is the only new step today. Everything else is Unit 2, unchanged.
  \end{itemize}
\end{frame}

% ── WHAT A PIECEWISE FUNCTION IS ──────────────────────────────────────────────
\begin{frame}[t, plain]
  \royalheader{One Function, Written in Pieces}
  \sectionlabel{\texttt{AFDA.AF.2} --- last night's delivery fee, now with a name}
  \begin{center}
  \begin{tikzpicture}
  \begin{axis}[
      width=9.2cm, height=3.4cm,
      xlabel={$m$ (miles)}, ylabel={$F$ (fee, \$)},
      xmin=0, xmax=8, ymin=0, ymax=16,
      xtick={0,1,2,3,4,5,6,7,8}, ytick={0,5,10,15}, minor y tick num=4,
      grid=both, grid style={linegray!45}, minor grid style={linegray!30},
      axis lines=left, tick label style={font=\tiny}, label style={font=\tiny}]
    \addplot[royal, very thick] coordinates {(0,5) (3,5) (8,15)};
  \end{axis}
  \end{tikzpicture}
  \end{center}
  \vspace{-0.35cm}
  \[
    F(m) = \begin{cases}
      5 & \text{if } 0 \le m \le 3,\\
      5 + 2(m-3) & \text{if } 3 < m \le 8.
    \end{cases}
  \]
  % A negative \vspace between a \[...\] display and the paragraph that follows it
  % collided the paragraph's first two lines (the display's depth is already
  % absorbed by \belowdisplayskip). Never pull text up under a display on a frame
  % -- shrink the figure instead, as done above.
  \small
  \textbf{Read the brace as ``if.''} $F(2)=5$ from the first piece; $F(6)=11$ from the second.
  \textbf{The pieces never overlap}, so every input still has exactly one output --- Lesson 2.1's
  definition of a function, unchanged.
\end{frame}

% ── CONTINUITY: CORNER VS BREAK ───────────────────────────────────────────────
\begin{frame}[t, plain]
  \royalheader{At the Boundary: Do the Two Rules Agree?}
  \sectionlabel{the paycheck --- and the answer settles the hook}
  \begin{center}
  \begin{tikzpicture}
  \begin{axis}[
      width=9.4cm, height=3.6cm,
      xlabel={$h$ (hours worked)}, ylabel={$P$ (pay, \$)},
      xmin=0, xmax=50, ymin=0, ymax=700,
      xtick={0,10,20,30,40,50}, ytick={0,240,480}, minor y tick num=1, minor x tick num=4,
      grid=both, grid style={linegray!45}, minor grid style={linegray!30},
      axis lines=left, tick label style={font=\tiny}, label style={font=\tiny}]
    \addplot[goldacc, thick, dashed] coordinates {(40,0) (40,700)};
    \addplot[royal, very thick] coordinates {(0,0) (40,480) (50,660)};
  \end{axis}
  \end{tikzpicture}
  \end{center}
  \vspace{-0.35cm}
  \small
  At $h=40$: \ $12(40)=480$, \ and the second rule starts at $480$ too. \textbf{They agree}, so the
  graph has no gap --- it is \textbf{continuous}, traceable without lifting the pencil.\\[0.1cm]
  \textbf{The graph turns a sharp corner at $(40,480)$ --- and a corner is \emph{not} a break.} The
  slope changed because the pay rate did. \quad And $P(45)=\$570$, which is student B.
\end{frame}

% ── THE JUMP AND THE DOTS ─────────────────────────────────────────────────────
\begin{frame}[t, plain]
  \royalheader{When the Pieces Disagree: a Jump, and Two Circles}
  \sectionlabel[goldacc]{shipping by weight --- \$6, then \$10, then \$16}
  \begin{center}
  \begin{tikzpicture}
  \begin{axis}[
      width=9.4cm, height=3.6cm,
      xlabel={$w$ (pounds)}, ylabel={$S$ (\$)},
      xmin=0, xmax=10.6, ymin=0, ymax=20,
      xtick={0,1,2,3,4,5,6,7,8,9,10}, ytick={0,6,10,16},
      grid=both, grid style={linegray!45},
      axis lines=left, tick label style={font=\tiny}, label style={font=\tiny}]
    \addplot[royal, very thick] coordinates {(0,6) (2,6)};
    \addplot[royal, very thick] coordinates {(2,10) (5,10)};
    \addplot[royal, very thick] coordinates {(5,16) (10,16)};
    \addplot[royal, only marks, mark=*, mark size=2pt] coordinates {(2,6) (5,10) (10,16)};
    \addplot[royal, only marks, mark=*, mark size=2pt, mark options={fill=white}]
      coordinates {(0,6) (2,10) (5,16)};
  \end{axis}
  \end{tikzpicture}
  \end{center}
  \vspace{-0.35cm}
  \small
  \textbf{What is $S(2)$?} Look straight up from $w=2$: one \textbf{filled} circle at \$6, one
  \textbf{open} circle at \$10. \quad \textbf{$S(2)=\$6$.}\\[0.1cm]
  \textbf{Exactly one filled circle sits above every input --- and it has to.} Two would mean two
  outputs for one input, and then it would not be a function at all.
\end{frame}

% ── THREE CONSEQUENCES ────────────────────────────────────────────────────────
\begin{frame}[t, plain]
  \royalheader{Three Things a Step Graph Does That Surprise People}
  \sectionlabel{all three read straight off the shipping picture}
  \begin{enumerate}
    \setlength{\itemsep}{0.28cm}
    \item \textbf{It breaks.} You lift the pencil at $w=2$ and $w=5$: \textbf{discontinuous}, with
          jumps of \$4 and \$6.
    \item \textbf{Its range is a list, not an interval:} $\{6, 10, 16\}$. No interval notation can
          say that --- exactly like the hot dogs in Lesson 2.1.
    \item \textbf{It is never increasing.} Every piece is \emph{constant}, and a piece is the only
          place ``increasing'' can be measured.
  \end{enumerate}
  \vspace{0.2cm}
  \textbf{The cost really does go up --- but it goes up \emph{at} the jumps, and a jump happens over
  no interval at all.}
\end{frame}

% ── THE WHOLE CHECKLIST STILL WORKS ───────────────────────────────────────────
\begin{frame}[t, plain]
  \royalheader{Every Characteristic in This Unit Still Works}
  \sectionlabel{\texttt{AFDA.AF.2h} --- one function, so you ask each question once}
  \renewcommand{\arraystretch}{1.3}
  \small
  \begin{tabularx}{\textwidth}{@{}p{4.2cm} p{4.4cm} X@{}}
    \textbf{Question} & \textbf{The paycheck $P$} & \textbf{The shipping cost $S$} \\
    \hline
    Domain, range (2.1) & $[0,50]$, \ $[0,660]$ & $(0,10]$, \ the list $\{6,10,16\}$ \\
    Zeros, intercept (2.4) & one zero at $h=0$; $(0,0)$ & no zeros at all \\
    Inc.\ / dec.\ / const.\ (2.4) & increasing on all of $[0,50]$ & constant on each piece; never
      increasing \\
    Absolute max, min (2.4) & \$660 at $h=50$; \$0 at $h=0$ & \$16 on $(5,10]$; \$6 on $(0,2]$ \\
    End behavior (2.5) & none --- the domain stops & none --- the domain stops \\
  \end{tabularx}
  \vspace{0.25cm}

  \textbf{Notice the two right-hand extremes are attained on a whole \emph{stretch}, not at one
  point.} An extreme is still a value \emph{and} a location --- the location just got wider.
\end{frame}

% ── GUIDED PRACTICE ───────────────────────────────────────────────────────────
\begin{frame}[t, plain]
  \royalheader{Guided Practice: The Delivery Drone}
  \sectionlabel{climb, cruise, descend --- three pieces, one flight}
  \begin{center}
  \begin{tikzpicture}
  \begin{axis}[
      width=9.6cm, height=3.6cm,
      xlabel={$t$ (seconds after take-off)}, ylabel={$A$ (feet)},
      xmin=0, xmax=14, ymin=0, ymax=60,
      xtick={0,2,4,6,8,10,12,14}, ytick={0,10,20,30,40,50}, minor x tick num=1,
      grid=both, grid style={linegray!45}, minor grid style={linegray!30},
      axis lines=left, tick label style={font=\tiny}, label style={font=\tiny}]
    \addplot[royal, very thick, domain=0:5, samples=60]{20*x-2*x^2};
    \addplot[royal, very thick] coordinates {(5,50) (9,50) (14,0)};
  \end{axis}
  \end{tikzpicture}
  \end{center}
  \vspace{-0.35cm}
  \small
  Increasing on $[0,5]$, constant on $[5,9]$, decreasing on $[9,14]$. \quad Zeros at $t=0$ and
  $t=14$ --- take-off and landing.\\[0.1cm]
  \textbf{Absolute maximum: $50$ feet --- and read its \emph{location} carefully.} The drone is at
  $50$ feet at every second from $t=5$ to $t=9$, so the maximum is attained on the whole interval
  $[5,9]$, not at one point.
\end{frame}

% ── ACTIVITY LAUNCH ───────────────────────────────────────────────────────────
\begin{frame}[t, plain]
  \royalheader{Group Activity: Which Piece Owns This Input?}
  \sectionlabel{groups of three --- everyone starts at Tier R}
  Every answer is still read off a picture.
  \begin{block}{The rule for the whole activity}
    Before you use a rule or read a value, \textbf{say out loud which piece owns the input}. And at
    a boundary, \textbf{let the dots decide} --- filled means included, open means not.
  \end{block}
  \vspace{0.1cm}
  Tier A closes on a graph with an \textbf{open circle at the top}, where the outputs get as close to
  $6$ as you like and never reach it --- so there is \emph{no absolute maximum}. That is Lesson 2.5's
  sentence in a completely new costume.
\end{frame}

% ── CLOSE ─────────────────────────────────────────────────────────────────────
\begin{frame}[t, plain]
  \royalheader{Exit Ticket}
  \sectionlabel[goldacc]{notes closed --- 5 minutes}
  \begin{enumerate}
    \setlength{\itemsep}{0.3cm}
    \item The babysitting rate: $B(3)$ and $B(6)$, which piece owns each, and the \textbf{boundary
          point}.
    \item Absolute maximum? Absolute minimum? \textbf{Read the minimum's location carefully} --- it
          is not a single point.
    \item Somebody wrote ``the graph has a break at $h=4$, because that is where the rule changes.''
          Naming $h=4$ was \emph{correct} in item 1. Say what is wrong anyway.
  \end{enumerate}
  \vspace{0.2cm}
  {\color{royal}\bfseries Next: Lesson 2.7 --- Choosing and Comparing Models in Context.}
\end{frame}

\end{document}
